\section{Refinement-Typed Calculus and Module Composition}
\label{sec:calculus}

We formalise \textsc{TensorGuard}'s abstract domain as a
refinement-type system over a \texttt{torch.fx}-style IR. The
type of a tensor is
\[
  \mathsf{Tensor}\bigl[\tau,\;\varphi_{\mathsf{shape}}(d_1,\dots,d_k),\;\varphi_{\mathsf{grad}}\bigr],
\]
where $\tau$ is a dtype, $\varphi_{\mathsf{shape}}$ is a
quantifier-free predicate over symbolic dim variables in
$\mathcal{T}_{\mathsf{shape}}=\mathrm{LIA}\cup\mathrm{Div}\cup\mathrm{BMul}$
(linear integer arithmetic, divisibility, and bounded
nonlinear multiplication restricted to syntactically-known
factors), and the grad refinement
$\varphi_{\mathsf{grad}}\in\{\mathsf{has\_grad},\mathsf{no\_grad},\top\}$
abstracts $\mathtt{requires\_grad}\wedge\mathtt{is\_on\_tape}$.
The grad lattice is flat: $\mathsf{has\_grad}$ and
$\mathsf{no\_grad}$ are incomparable, joined to $\top$. A
typing context $\Gamma$ binds variables to refinement types,
accumulates path predicates, and tracks a grad-context flag
$G\in\{\mathsf{on},\mathsf{off}\}$ recording whether the cursor
is inside \texttt{torch.no\_grad()}. The entailment
$\Gamma\models\varphi$ is discharged by Z3 on the shape
projection and by the lattice on the grad component.

\paragraph{Representative typing rules.}
We give five rules; the full set (matmul, conv1d/2d/3d,
broadcast, view, reshape, split, chunk, unbind, cat, stack,
permute, transpose, einsum, indexing, slicing, scatter,
gather, repeat, expand, SDPA, LayerNorm, RMSNorm,
\texttt{detach}, \texttt{no\_grad} context, optimizer step) is in
\Cref{app:typing-rules}, with pen-and-paper soundness in
\Cref{app:soundness}.

{\small
\begin{mathpar}
\inferrule*[lab=\textsc{T-Matmul}]
{\Gamma\vdash a:\mathsf{Tensor}[\tau,(\bar B,m,k),\varphi_a]\\
 \Gamma\vdash b:\mathsf{Tensor}[\tau,(\bar B',k,n),\varphi_b]\\
 \Gamma\models\mathrm{broadcast}(\bar B,\bar B')=\bar C}
{\Gamma\vdash a@b:\mathsf{Tensor}[\tau,(\bar C,m,n),g(a,b)]}
\and
\inferrule*[lab=\textsc{T-View($-1$)}]
{\Gamma\vdash x:\mathsf{Tensor}[\tau,(d_1,\dots,d_n),\varphi_x]\\
 \bar s=(s_1,\dots,s_{j-1},-1,s_{j+1},\dots,s_m)\\
 P=\textstyle\prod d_i,\;\;Q=\textstyle\prod_{i\neq j}s_i\\
 \Gamma\models Q\mid P}
{\Gamma\vdash x.\mathrm{view}(\bar s):\mathsf{Tensor}[\tau,(s_1,\dots,P/Q,\dots,s_m),\varphi_x]}
\and
\inferrule*[lab=\textsc{T-Detach}]
{\Gamma\vdash x:\mathsf{Tensor}[\tau,\varphi_s,\varphi_g]}
{\Gamma\vdash x.\mathrm{detach}():\mathsf{Tensor}[\tau,\varphi_s,\mathsf{no\_grad}]}
\end{mathpar}}

\paragraph{Verdicts.} The verifier emits one of five verdicts:
\textsc{Verified}, \textsc{Refuted-Proof} (RP, an unconditional
witness), \textsc{Contract-Violation} (CV, sound under the
caller-rely $\mathit{assume}_e$), \textsc{Library-Warn} (LW, a
conservative warning on dispatch outside the modeled fragment),
and \textsc{Abstain}. The soundness theorem
(\Cref{thm:soundness}) covers RP and CV; LW and \textsc{Abstain}
make no claim.

%% fragment_v8.tex — modeled fragment definition and fragment-level soundness
\subsection{The Modeled Fragment $\mathcal{F}_\mathrm{TG}$}
\label{sec:fragment}

\noindent\textbf{Syntax.} The fragment $\mathcal{F}_\mathrm{TG}$ is defined by the following grammar,
where $\mathit{op}\in\mathrm{Cat}$ denotes any operator in the catalogue of
Table~\ref{tab:handler-soundness}, and static integers are values derivable
from the declared refinements without a runtime query.

{\small
\[
\begin{array}{@{}lll@{}}
M &::=& \textbf{class}\ C(\mathtt{nn.Module}){:}\ \textbf{\_\_init\_\_}(\ldots){;}\ \textbf{forward}(\mathit{self}, x_1,\dots,x_n) \to s \\[2pt]
s &::=& x = e \mid \textbf{return}\ e \mid \textbf{if}^{\dagger}\ c{:}\ s_1\ \textbf{else}\ s_2
 \mid \textbf{for}^{\dagger}\ x\ \textbf{in}\ \textbf{range}(n){:}\ s \\[2pt]
e &::=& x \mid e.\mathit{op}(\mathit{args})\ [\mathit{op}\in\mathrm{Cat}]
 \mid C(\mathit{args}) \mid \mathit{self}.\mathit{attr} \mid \mathit{const}
 \mid (e_1,\dots,e_k) \mid x_1,\dots,x_k = e
\end{array}
\]
}

\noindent${}^\dagger$\emph{Shape-determinism restriction}: branch condition $c$ and loop
bound $n$ must be static integers derivable from the declared refinements.
Types are $\mathsf{Tensor}[\tau,\varphi_s,\varphi_g]$ as in \S\ref{sec:calculus}.

\smallskip\noindent\textbf{Out-of-fragment.}
The following constructs are \emph{not} in $\mathcal{F}_\mathrm{TG}$:
data-dependent control flow (condition involves a runtime tensor value);
custom \texttt{autograd.Function}; dynamic attribute mutation of a module
after \texttt{\_\_init\_\_}; \texttt{exec}/\texttt{eval}; decorator
metaprogramming that modifies \texttt{forward} at runtime.
The analyser emits \textsc{Abstain} on any block containing these.

\begin{theorem}[Fragment-level soundness on $\mathrm{Cat}_{\mathrm{sound}}$]
\label{thm:fragment-soundness}
Let $M\in\mathcal{F}_\mathrm{TG}$ be a closed module \emph{such that
every operator occurrence in $M$ uses a handler in
$\mathrm{Cat}_{\mathrm{sound}}=\mathrm{Cat}_{\mathrm{audit}}\cup
\mathrm{Cat}_{\mathrm{pen}}$}. Let $\Gamma\vdash M:\tau$ with
$\mathsf{Verify}(M)=\textsc{Verified}$. Then for every input $x$
with $\sigma\models\mathit{assume}_M(x)$, the reduction
$\langle M(x),\sigma\rangle\to^* v$ does not invoke the
operator-level shape-mismatch exception of any operator in
$\mathrm{Cat}_{\mathrm{sound}}$.
\end{theorem}
\begin{proof}[Sketch]
\emph{Preservation}: by induction on the BNF of $\mathcal{F}_\mathrm{TG}$.
Each operator rule in $\mathrm{Cat}_{\mathrm{sound}}$
(\textsc{T-Matmul}, \textsc{T-View($-1$)}, \textsc{T-Reshape($-1$)},
\textsc{T-Detach}, and the remaining $40$ rules in
\Cref{app:typing-rules} restricted to $\mathrm{Cat}_{\mathrm{sound}}$)
produces a uniquely-typed result. The $28$ Lean-audited handlers
(Table~\ref{tab:handler-soundness}, top block) are sorry-free at
the rule definitions; the $16$ pen-and-paper handlers follow from
the operator-level argument in \Cref{app:soundness}, and two of
their nontrivial integer-arithmetic obligations
(\texttt{reshape} divisibility and the
single-inferred-dim case) are also mechanised in Lean as
\texttt{reshape\_sound\_zero\_unknowns},
\texttt{reshape\_sound\_one\_unknown}, and
\texttt{reshape\_rejects\_multi\_unknown} (\Cref{sec:eval-lean}).
\emph{Progress}: each statement in $\mathcal{F}_\mathrm{TG}$ either reduces
or is a value; the shape-determinism restriction on $\textbf{if}^\dagger$ and
$\textbf{for}^\dagger$ ensures branch selection and loop unrolling are
determined by the static refinements, so no stuck state arises at those
forms.
Together, no well-typed term using only operators in
$\mathrm{Cat}_{\mathrm{sound}}$ raises a shape-mismatch at any
$\mathit{op}\in\mathrm{Cat}_{\mathrm{sound}}$. Coverage of the
$35$ tested-only handlers
$\mathrm{Cat}_{\mathrm{tested}}=\mathrm{Cat}\setminus
\mathrm{Cat}_{\mathrm{sound}}$ is the subject of
\Cref{conj:tested-only-soundness}, not of this theorem.
\end{proof}

\paragraph{Boundary precision.}
Theorem~\ref{thm:fragment-soundness} does \emph{not} claim:
(a)~termination;
(b)~numerical correctness or absence of NaN/Inf;
(c)~absence of dtype errors outside $\mathrm{Cat}_{\mathrm{sound}}$;
(d)~soundness of the analyser \emph{implementation}---only of the rule
table it implements;
(e)~coverage of the $35$ tested-only handlers
$\mathrm{Cat}_{\mathrm{tested}}$ (Table~\ref{tab:handler-soundness},
addressed by \Cref{conj:tested-only-soundness}).
The implication ``verdict = \textsc{Verified} $\Rightarrow$ no shape mismatch''
depends on: (i)~the $\mathit{assume}_M$ contract holding at the call site;
(ii)~the operator being in $\mathrm{Cat}_{\mathrm{sound}}$ with a
Lean-audited or pen-and-paper soundness proof; (iii)~Z3 returning
\textsc{Unsat} correctly on every discharged obligation.



\noindent The fragment-level statement (Thm.~\ref{thm:fragment-soundness})
specialises Thm.~\ref{thm:soundness} to whole programs.

\begin{theorem}[Soundness of the proof-grade verdicts]
\label{thm:soundness}
For any closed $e$ in the supported fragment and any heap
$\sigma\models\Gamma$:
(i) if $\mathsf{Verify}(e)=\textsc{Verified}$, no reduction
sequence from $\langle e,\sigma,\mathsf{on}\rangle$ raises a
shape-mismatch exception, and the grad-flag claim holds at
every \texttt{loss.backward()} site;
(ii) if $\mathsf{Verify}(e)=\textsc{Refuted-Proof}$, every
reduction sequence raises a shape- or grad-flag exception at
the witnessed subterm;
(iii) if $\mathsf{Verify}(e)=\textsc{Contract-Violation}$, (ii)
holds under the assumption that every caller supplies inputs
satisfying $\mathit{assume}_e$;
(iv) if $\mathsf{Verify}(e)\in\{\textsc{Library-Warn},\textsc{Abstain}\}$,
no claim is made.
\end{theorem}
The proof (induction on reductions; full case analysis in
\Cref{app:soundness}) reduces (i),(ii) to per-operator
preservation lemmas already covered by the Lean rule audit
(\Cref{sec:eval-lean}), and (iii) to the same argument inside
the rely envelope $\mathit{assume}_e$.

\smallskip\noindent\textbf{Subject reduction and progress.}
The classical Preservation and Progress theorems for $\mathcal{F}_{\mathrm{TG}}$
are established in \Cref{app:subject-reduction}
(\Cref{thm:preservation,thm:progress}).
Together they yield \Cref{cor:soundness}: well-typed programs in
$\mathcal{F}_{\mathrm{TG}}$ whose operator preconditions hold in $\sigma$ raise no
shape-mismatch exception on any reduction sequence.
Theorem~\ref{thm:soundness} is the abstract-interpretation specialisation of this
classical pair to the verdict lattice; the verdict \textsc{Verified} corresponds
to Z3 discharging all preconditions symbolically.

\subsection{Assume/Guarantee Composition}
\label{sec:assume-guarantee}

For each \texttt{nn.Module} class $M$ we synthesise a contract
$C(M)=(\mathit{assume}_M,\mathit{guarantee}_M)$, with
$\mathit{assume}_M$ a refinement on input types and
$\mathit{guarantee}_M$ a refinement on output types. Sequential
composition (e.g.\ \texttt{nn.Sequential}) discharges
$\mathit{guarantee}_{M_1}\!\Rightarrow\!\mathit{assume}_{M_2}$
locally; parallel composition (residual sums, QKV split) treats
each sibling independently and joins the guarantees:

{\small
\begin{mathpar}
\inferrule[\textsc{Seq}]
{\vdash C(M_1)=(A_1,G_1)\\
 \vdash C(M_2)=(A_2,G_2)\\
 \models G_1\Rightarrow A_2}
{\vdash C(M_2\circ M_1)=(A_1,G_2)}
\and
\inferrule[\textsc{Sub}]
{C(M)=(A,G)\\ M'<:M\\
 \models A_M\Rightarrow A_{M'}\quad\text{(contravariant)}\\
 \models G_{M'}\Rightarrow G_M\quad\text{(covariant)}}
{\vdash C(M')=(A_{M'},G_{M'})\text{ refines }C(M)}
\end{mathpar}}

\noindent The subclassing rule is a syntactic obligation: a
subclass that overrides \texttt{forward} in a way that violates
the parent's guarantee is rejected statically rather than
silently re-analysed.

\begin{theorem}[Compositional soundness; mechanised fragment]\label{thm:ag-sound}
\emph{Scope.} The theorem as stated is mechanised in
\texttt{lean/TensorGuard/AssumeGuarantee.lean} on the same
$3$-operator DSL (\texttt{matmul}/\texttt{view}/\texttt{add})
as \texttt{Soundness.lean}; it does \emph{not} cover the full
$79$-handler operator surface of \Cref{tab:handler-soundness},
the AST extractor, or the Python analyser implementation.
If every leaf module in the call tree of an
$\mathcal{F}_{\mathrm{TG}}^{3}$ program is contract-verified
in the mechanised fragment and every composition obligation
is discharged by Z3, then the whole-program verdict given by
the inductive use of \textsc{Seq}, \textsc{Par}, and
\textsc{Sub} is sound in the sense of \Cref{thm:soundness}
restricted to the same fragment.
\end{theorem}

\begin{theorem}[Monotonicity of refinement]\label{thm:monotonicity}
Strengthening any single contract $C(M_0)$ to $C'(M_0)\sqsubseteq
C(M_0)$ never turns a \textsc{Verified} root into
\textsc{Refuted-Proof}; the new verdict is in
$\{\textsc{Verified},\textsc{Contract-Violation},
\textsc{Library-Warn},\textsc{Abstain}\}$.
\end{theorem}

Proofs (and the rely/guarantee axiom of fresh refutation
witnesses needed to make \Cref{thm:monotonicity} hold) are in
\Cref{app:ag-proofs}.
